\documentclass{article}

\usepackage{amsmath,amssymb}
\usepackage{xurl}
\usepackage[hidelinks]{hyperref}
\setlength{\emergencystretch}{2em}

\input{command}

\title{The General Compact-Convex Brouwer Fixed-Point Theorem}
\author{}
\date{2026-08-04}

\begin{document}
\maketitle
\gapnote{open-gap}{resolved}

\section{The assertion}

Wolf's Theorem~6.10 states the Brouwer fixed-point theorem in its general
compact-convex form:\footnote{\cite[Theorem~6.10]{Wolf2012Quantum}; local source
\path{Notes/WolfNoteTexSource/ch06\_spectral\_properties.tex},
lines~1143--1151.}
\begin{align}
  T : S \to S,\qquad
  S \subset \mathbb{R}^n \text{ nonempty, compact, convex},\qquad
  T \text{ continuous}
  \ \Longrightarrow\
  \exists\, x \in S,\ T(x) = x.
\end{align}

\section{Resolution}

The vendored Gametheory library (LionSR/Brouwer) proves the Brouwer fixed-point
theorem for the standard simplex~$\Delta^n$:
\leanid{Brouwer}.
QICLean extends this to products of simplices, closed cubes, and compact retracts
of finite-dimensional real normed spaces.\footnote{Formal declarations:
\leanid{exists\_fixedPoint\_closedCube} in
\doclink{QICLean/Topology/BrouwerProduct},
\leanid{fixedPoint\_of\_compact\_retract} in
\doclink{QICLean/Topology/CompactRetractFixedPoint},
and \leanid{fixedPoint\_of\_compact\_convex} and
\leanid{brouwer\_fixedPoint\_compactConvex} in
\doclink{QICLean/Topology/CompactConvexFixedPoint}.}

The last two declarations now prove the coordinate-free theorem and Wolf's
printed form for an arbitrary nonempty compact convex subset
\(S \subset \mathbb{R}^n\), respectively.  The proof follows the metric-projection
route previously identified in this note.  For each ambient point, compactness
provides a minimizer of squared distance in~\(S\).  Convexity and the
first-order variational inequality imply that the chosen projection is
1-Lipschitz and restricts to the identity on~\(S\).  It therefore supplies the
continuous retraction required by
\leanid{fixedPoint\_of\_compact\_retract}.  The reusable projection declarations
are \leanid{CompactConvex.metricProjection},
\leanid{CompactConvex.nearest\_inner\_nonpos}, and
\leanid{CompactConvex.metricProjection\_lipschitzWith}.

Wolf supplies only the theorem statement at the cited location, so this route
is recorded as the formal proof and is not attributed to the lecture notes.

\section{Impact}

The former general-statement gap did \emph{not} block the formalization of Wolf
Theorem~6.11 (Stationary states), Proposition~6.8 (Positive fixed-points), or
Corollary~6.5 (spanning by stationary density matrices).  These results rely
only on the density-matrix version of Brouwer, which is already proved.
The density-matrix specialization remains the theorem used by those channel
results; the general theorem is now available independently.

\section{Verdict}

\textbf{Resolved.}  Wolf's general compact-convex statement is formalized as
\leanid{brouwer\_fixedPoint\_compactConvex}.  The coordinate-free theorem and
the reusable nonexpansive metric projection are formalized alongside it.

\bibliographystyle{alpha}
\bibliography{references}

\end{document}
